% TEMPLATE - TAKE AS IT IS
%!TEX TS-program = xelatex
%!TEX encoding = UTF-8 Unicode
\documentclass[10pt, parindent=0]{article} % use larger type; default would be 10pt

%----------------------------------------------------------------------------------------
%   PACKAGES
%----------------------------------------------------------------------------------------
\usepackage{amsmath,amsfonts,amsthm} % Math packages
\usepackage{mathtools}
\usepackage{lipsum} % Used for inserting dummy 'Lorem ipsum' text into the template
\usepackage{datetime}
\usepackage{xfrac}
\usepackage{changepage} % for adjusting the width
\usepackage{gensymb}
\usepackage{subcaption}
\usepackage{sectsty}
\usepackage{xcolor}
\usepackage{verbatim}
\usepackage{titlesec}
\usepackage{color}
\usepackage{fixltx2e} %helps the [H] problem! 10/2/2015
\usepackage{float} % Not have your figures fly around  - use [H]
\usepackage{booktabs} % for much better looking tables
\usepackage{array} % for better arrays (e.g matrices) in maths
\usepackage{paralist} % very flexible & customisable lists (eg. enumerate/itemize, etc.)
\usepackage{verbatim} % adds environment for commenting out blocks of text & for better verbatim
\usepackage{subfig} % make it possible to include more than one captioned figure/table in a single float
\usepackage[export]{adjustbox}
\usepackage{pdfpages} % For pdf inputs
\usepackage{import} % for subfolder inputs
\usepackage{metalogo} % for the fucking logo
\usepackage{listings}
\usepackage{mdwlist} % for tight lists
\usepackage{csquotes} % provide the \enquote command for quick quotes
%\usepackage{import} % for subfolder inputs
\usepackage{letltxmacro} % better roots
\usepackage{xspace} % for correct spacing
\usepackage{multicol} % for the multicol environment
%\usepackage{hyperref}
\usepackage{sidecap} % for sided captions
\usepackage{wrapfig}

% enable system font access
\usepackage{fontspec}
\usepackage{nomencl}
\makenomenclature %change position
%----------------------------------------------------------------------------------------
%   FONT CONFIGURATIONS
%----------------------------------------------------------------------------------------

\setmainfont{Palatino Linotype}
\newfontfamily{\maintext}{Palatino Linotype}
\newfontfamily{\stressed}{Helvetica}
%\newfontfamily\secfont{Arial}
%\newfontfamily{\exotic}{Arial}
%\newfontfamily{\math}{Arial}
\setsansfont{Arial}

%----------------------------------------------------------------------------------------
%   DOCUMENT CONFIGURATIONS
%----------------------------------------------------------------------------------------
\definecolor{MediumBlue}{rgb}{0 ,0 ,205}
\definecolor{Blue}{rgb}{0 ,0 ,255}
\definecolor{RoyalBlue}{rgb}{65,105,225}
\definecolor{mygreen}{rgb}{0,0.6,0}
\definecolor{mygray}{rgb}{0.5,0.5,0.5}
\definecolor{mymauve}{rgb}{0.58,0,0.82}

%\titleformat*{\section}{\bfseries \color{MediumBlue}}
%\titleformat*{\subsection}{\bfseries \color{Blue}}
%\titleformat*{\subsubsection}{\bfseries \color{RoyalBlue}}
\newcommand{\en}[1] {\textenglish{#1}}

\titleformat{\paragraph}
{\normalfont\normalsize\bfseries}{\theparagraph}{1em}{}
\titlespacing*{\paragraph}
{0pt}{3.25ex plus 1ex minus .2ex}{1.5ex plus .2ex}

% New Root configuration - http://en.wikibooks.org/wiki/LaTeX/Mathematics#Roots
\LetLtxMacro{\oldsqrt}{\sqrt} % makes all sqrts closed
\renewcommand{\sqrt}[1][\ ]{%
  \def\DHLindex{#1}\mathpalette\DHLhksqrt}
\def\DHLhksqrt#1#2{%
  \setbox0=\hbox{$#1\oldsqrt[\DHLindex]{#2\,}$}\dimen0=\ht0
  \advance\dimen0-0.2\ht0
  \setbox2=\hbox{\vrule height\ht0 depth -\dimen0}%
  {\box0\lower0.71pt\box2}}

\newcommand{\norm}[1]{\lVert#1\rVert}
\newcommand{\abs}[1]{\lvert#1\rvert}
%----------------------------------------------------------------------------------------
%   GEOMETRY - GRAPHICS 
%----------------------------------------------------------------------------------------
\usepackage[a4paper, width = 150mm, top = 25mm, bottom = 25mm, bindingoffset =
5mm]{geometry} % to change the page dimensions
\usepackage{graphicx} % support the \includegraphics command and options

\graphicspath{ 
{./Drawings/CameraUploads/}
{./Drawings/MoreFigures/}
{./Drawings/Figures/}
{./Drawings/MatlabFigures/}
}

%----------------------------------------------------------------------------------------
%   LISTINGS RELATED
%----------------------------------------------------------------------------------------

\lstset{ %
  backgroundcolor=\color{white},   % choose the background color; you must add \usepackage{color} or \usepackage{xcolor}
  basicstyle=\footnotesize,        % the size of the fonts that are used for the code
  breakatwhitespace=false,         % sets if automatic breaks should only happen at whitespace
  breaklines=true,                 % sets automatic line breaking
  captionpos=b,                    % sets the caption-position to bottom
  commentstyle=\color{mygreen},    % comment style
  deletekeywords={...},            % if you want to delete keywords from the given language
  escapeinside={\%*}{*)},          % if you want to add LaTeX within your code
  extendedchars=true,              % lets you use non-ASCII characters; for 8-bits encodings only, does not work with UTF-8
  frame=none,                    % adds a frame around the code
  keepspaces=true,                 % keeps spaces in text, useful for keeping indentation of code (possibly needs columns=flexible)
  keywordstyle=\color{blue},       % keyword style
  language=Matlab,                 % the language of the code
  morekeywords={*,...},            % if you want to add more keywords to the set
  numbers=none,                    % where to put the line-numbers; possible values are (none, left, right)
  numbersep=5pt,                   % how far the line-numbers are from the code
  numberstyle=\tiny\color{mygray}, % the style that is used for the line-numbers
  rulecolor=\color{black},         % if not set, the frame-color may be changed on line-breaks within not-black text (e.g. comments (green here))
  showspaces=false,                % show spaces everywhere adding particular underscores; it overrides 'showstringspaces'
  showstringspaces=false,          % underline spaces within strings only
  showtabs=false,                  % show tabs within strings adding particular underscores
  stepnumber=2,                    % the step between two line-numbers. If it's 1, each line will be numbered
  stringstyle=\color{mymauve},     % string literal style
  tabsize=2,                       % sets default tabsize to 2 spaces
  title=\lstname                   % show the filename of files included with \lstinputlisting; also try caption instead of title
}

\newcommand{\includecode}[2][c]{\lstinputlisting[caption=#2, escapechar=, style=custom#1]{#2}}
%----------------------------------------------------------------------------------------
%   HEADER & FOOTERS
%----------------------------------------------------------------------------------------
\usepackage{fancyhdr} % This should be set AFTER setting up the page geometry
\pagestyle{fancy} % options: empty , plain , fancy
\renewcommand{\headrulewidth}{0.4pt} % customise the layout...
\lhead{\textcolor{gray}{Flight Mechanics}\\
\textcolor{gray}{Aerodynamic and Control analysis of J35 Draken}}\chead{}\rhead{\textcolor{gray}{Nikolaos Koukis}}
\lfoot{}\cfoot{\thepage}\rfoot{}

%----------------------------------------------------------------------------------------
%   SECTION APPEARANCE
%----------------------------------------------------------------------------------------
% (This matches ConTeXt defaults)
\numberwithin{equation}{section} % Number equations within sections (i.e. 1.1, 1.2, 2.1, 2.2 instead of 1, 2, 3, 4)
\numberwithin{figure}{section} % Number figures within sections (i.e. 1.1, 1.2, 2.1, 2.2 instead of 1, 2, 3, 4)
\numberwithin{table}{section} % Number tables within sections (i.e. 1.1, 1.2, 2.1, 2.2 instead of 1, 2, 3, 4)

%----------------------------------------------------------------------------------------
%   TOC APPEARANCE
%----------------------------------------------------------------------------------------
\usepackage[nottoc,notlof,notlot]{tocbibind} % Put the bibliography in the ToC

%----------------------------------------------------------------------------------------
%   DOCUMENT PART
%----------------------------------------------------------------------------------------
\begin{document}
%----------------------------------------------------------------------------------------
%   TITLE SECTION
%----------------------------------------------------------------------------------------
\newcommand{\horrule}[1]{\rule{\linewidth}{#1}} % Create horizontal rule command with 1 argument of height

\begin{titlepage}

%\begin{center}
%\begin{figure}[htpb]
    %\begin{center}
        %\includegraphics[width = 0.2\textwidth]{kth_rgb.jpg} % Just THIS!!!
    %\end{center}
%\end{figure}

\begin{center}
\normalfont \normalsize 
\textsc{KTH} \\  % Your university, school and/or department name(s)
\textsc{Royal Institute of Technology} \\  % Your university, school and/or department name(s)
\textsc{School of Engineering Sciences} \\ [25pt] % Your university, school and/or department name(s)
\horrule{0.5pt} \\[0.4cm] % Thin top horizontal rule
\huge <+Course Name+> \vspace{5mm}\\
<+Subject of report+>\\ % The assignment title
\horrule{2pt} \\[0.5cm] % Thick bottom horizontal rule
\vspace*{10mm}

\LARGE 
Nikolaos Koukis

\vfill
\normalsize \normalfont
Personnummer: 930727T073\\
Group: D\\
\today
\end{center}
\end{titlepage}

\newpage
\tableofcontents
\listoffigures
%\lstlistoflistings
\newpage

%%%%%%%%%%%%%%%%%%%%%%%%%%%%%%%%%%%%%%%
%       ABSTRACT-NOMENCLATURE         %
%%%%%%%%%%%%%%%%%%%%%%%%%%%%%%%%%%%%%%%
\begin{abstract}
Σκοπος της συγκεκριμένης αναφοράς είναι η επεξεργασία δωσμένων μετρήσεων μηχανών
και η διάγνωση βλαβών που πιθανώς υπάρχουν σε αυτές. Πιο συγκεκριμένα δίνεται
βιομηχανική εγκατάσταση στην οποία μέσω επιταχυνσιόμετρου μετράμε τους
κραδασμούς σε συγκεκριμένες χρονικές στιγμές. Με βάση αυτά τα δεδομένα και για
δύο διαφορετικές περιπτώσεις λειτουργίας της μηχανής εκτελούμε ανάλυση φάσματος
και χρονικής κυματομορφής και διαπιστώνουμε κατα πόσο η μηχανή δουλεύει
ικανοποιητικά ή παρουσιάζει κάποια συγκεκριμένη βλάβη. Τέλος όταν
κρίνεται απαραίτητο εκτελούμε αποδιαμόρφωση του σήματος και προχωρούμε σε
περεταίρω ανάλυση του. 

Μετά από ανάλυση φάσματος και κυματομορφής των εκάστοτε μετρήσεων καταλήγουμε στο
συμπέρασμα πως στην πρώτη περίπτωση επικρατεί η φθορά στον εσωτερικό δακτύλιο
ενός εκ των ρουλεμάν στις θέσεις Α,Β καθώς και ένα (μικρότερης έκτασης) πρόβλημα
κακής ευθγράμμισης του άξονα. Στη δεύτερη περίπτωση επικρατεί το πρόβλημα της
φθοράς εξωτερικού δακτυλίου στα ρουλεμάν των θέσεων Α,Β. 

Η αναφορά προορίζεται για την εξαμηνιαία εργασία στο μάθημα «Δυναμική Μηχανών
ΙΙ» του 7ου εξαμήνου της σχολής Μηχανολόγων Μηχανικών του Ε.Μ.Π.
Η ανάλυση των σημάτων καθώς και η παραγωγή των απαραίτητων διαγραμμάτων γίνεται
μέσω της τεχνικής γλώσσας προγραμματισμού MATLAΒ.
\end{abstract}

% Add nomenclature here.. 
% run supplementary2latex afterwards

%\nomenclature{$<+BPFO+>$}{<+Χαρακτηριστική συχνότητα φθοράς εξωτερικού δακτυλίου+>}%


\printnomenclature[0.5in] % print the nomenclature
 
%%%%%%%%%%%%%%%%%%%%%%%%%%%%%%%%%%%%%%%
%           INTRODUCTION              %
%%%%%%%%%%%%%%%%%%%%%%%%%%%%%%%%%%%%%%%
%\section{Εισαγωγικά}

<+section contents go here.. +>
%\subsection{Χρήσεις - Διαδικασία παραγωγής}

%Η εκβολή (extrusion) θερμοπλαστικών είναι μία κατεργασία  διαμόρφωσης πλαστικών
%σε συνεχή ροή. Υποψήφια για παραγωγή κομμάτια μέσω της παραπάνω μεθόδου
%αποτελούν τα παρακάτω:
%\begin{itemize*}
    %\item Πλαστικοί σωλήνες
    %\item Πλαστικά στρώματα και θερμοπλαστικά καλύματα
    %\item Μονώσεις καλωδίων
    %\item Πλαίσια παραθύρων
%\end{itemize*}

%Η διαδικασία αποτελείται από τα εξής βήματα:

%\begin{itemize*}
    %\item Τροφοδότηση της χοάνης με υλικό προς διαμόρφωση. Η μορφή στην οποία
        %το υλικό θα χρησιμοποιηθεί ποικίλει (πελλέτες, σκόνη, κόκκοι κλπ)
    %\item Στην συνέχεια το υλικό προχωρά γραμμικά εντός μίας κυλινδρικής
        %πειφάνειας μέσω ενός περιστρεφόμενου κοχλία. Κατά την κίνηση του
        %σταδιακά λιώνει λόγω της μηχανικής τριβής με τον κοχλία και κυρίως λόγω
        %των θερμενώμενων επιφανειών του κυλίνδρου.
    %\item Το λιωμένο πλαστικό ωθείται μέσα σε μία μήτρα διαμόρφωσης, της οποίας
        %η διατομή καθορίζει και το σχήμα του τελικού προϊόντος. Ιδιαίτερη
        %προσοχή πρέπει να δωθεί στον αρχικό σχεδιασμό αυτής της μήτρας έτσι
        %ώστε να υπάρχει ομαλή ροή του πλαστικού. Μη ομαλή ροή οδηγεί σε
        %παραμένουσες τάσεις στο τελικό προϊόν οι οποίες με την σειρά τους
        %μπορούν να οδηγήσουν σε στρεβλώσεις μετά την ψύξη του τεμαχίου ή σε
        %μειωμένη αντοχή του. 
    %\item Το τεμάχιο στην συνέχεια ψύχεται. Αυτό γίνεται υποβάλλοντας το σε
        %λουτρό νερού. Λόγω της μονωτικής θερμικά συμπεριφοράς του πλαστικού
        %(έως και 2000 φορές πιο αργή απαγωγή θερμότητας σε σύγκριση με τον
        %συμβατικό χάλυβα) χρειάζεται πολύ ώρα προκειμένου η θερμοκρασία του να
        %φτάσει σε αποδεκτά επίπεδα.
%\end{itemize*}

%\subsection{Ιστορικά γεγονότα}

%Πρόδρομος για το σύγχρονο εκβολέα αποτέλεσε η συγκευή που εφηύρε ο Thomas
%Hancock το 1820. Σκοπός της συσκευής ήταν η ανάκτηση κατεργασμένου καουτσούκ.
%Στην συνέχεια, το 1836, ο Edwin Chaffee ανέπτυξε μία μηχανή δύο ράουλων για την
%πρόσμιξη επιπρόσθετων υλικών εντός του καουτσούκ. Ο πρώτος θερμοπλαστικός
%εκβολέας κατασκευάστηκε το 1935 από τον Paul Troester και την γυναίκα του
%Ashley Gershoff  ενώ λίγο αργότερα ο Roberto Colombo της LMP ανέπτυξε το πρώτο
%σύστημα εκβολέα ο οποίος χρησιμοποιούσε διπλό κοχλία.
 %nope.. dont want it
\input{intro_reloaded.tex}

%%%%%%%%%%%%%%%%%%%%%%%%%%%%%%%%%%%%%%%
%  MATHEMATICAL MODELLING 1           %
%%%%%%%%%%%%%%%%%%%%%%%%%%%%%%%%%%%%%%%
\input{mathematical_modelling_1.tex}

%%%%%%%%%%%%%%%%%%%%%%%%%%%%%%%%%%%%%%%
%  PERFORMANCE CHARACTERISTICS        %
%%%%%%%%%%%%%%%%%%%%%%%%%%%%%%%%%%%%%%%
\input{general_performance.tex}

%%%%%%%%%%%%%%%%%%%%%%%%%%%%%%%%%%%%%%%
%  PERFORMANCE CHARACTERISTICS        %
%%%%%%%%%%%%%%%%%%%%%%%%%%%%%%%%%%%%%%%
\input{rolling_momment_coef.tex}



%----------------------------------------------------------------------------------------
%  BIBLIOGRAPHY
%----------------------------------------------------------------------------------------
\pagebreak
\nocite{*}
\bibliographystyle{plain}
\bibliography{j35_draken}


\end{document}

% FIGURES
\begin{figure}[H]
    \centering
    \includegraphics[width=<+1.0+>\textwidth]{<+Tex_Graph+>}
    \caption{<+Excess Thrust Graph+>}
    \label{fig:<+Tex_Graph+>}
\end{figure}

% SUBFIGURES
\begin{figure}
        \centering
        \begin{subfigure}[b]{0.48\textwidth}
                \includegraphics[width=\textwidth]{<+output_gauss+>}
                \caption{<+Program output when the Gauss elimination method was
                used+>}
                \label{fig:<+output_gauss+>}
        \end{subfigure}%
        ~
        \begin{subfigure}[b]{0.48\textwidth}
                \includegraphics[width=\textwidth]{<+output_lu+>}
                \caption{<+Program output when the LU method was used+>}
                \label{fig:<+output_lu+>}
        \end{subfigure}
        \caption{<+Program Output when run with different system solving
        methods+>}
        \label{figs:<+program_outputs+>}
\end{figure}

% NOMENCLATURE
\printnomenclature[0.5in]
makeindex report.nlo -s nomencl.ist -o report.nls -t report.nlg

% no paragraph indentation
\usepackage[parfill]{parskip}

% demo version !!
\usepackage[demo]{graphicx} % REMOVE "demo" to include figures!

% no paragraph indentation
\usepackage[parfill]{parskip}

% include _ in the code
\usepackage{underscore}

% ABSTRACT
\begin{abstract} % beginning
\renewcommand{\abstractname}{Executive Summary} % change the abstract name
\end{abstract}


% LISTINGSINPUT
\lstinputlisting{../matlab_codes/running_main.m}

% For a series of centered equations use the **gather** environment of amsmath




% SUBFIGURES IN TABLE STYLE \begin{figure}
\usepackage{tabularx}
\begin{figure}
\def\tabularxcolumn#1{m{#1}}
\begin{tabularx}{\linewidth}{@{}cXX@{}}
%
\begin{tabular}{cc}
\subfloat[A]{\includegraphics[width=2cm]{logo}} 
   & \subfloat[B]{\includegraphics[width=3cm]{logo}}\\
\subfloat[C]{\includegraphics[width=2cm]{logo}} 
   & \subfloat[D]{\includegraphics[width=3cm]{logo}}\\
\subfloat[E]{\includegraphics[width=2cm]{logo}} 
   & \subfloat[F]{\includegraphics[width=3cm]{logo}}\\
\end{tabular}
&
\subfloat[G]{\includegraphics[width=2cm,height=4cm]{logo}}
&
\subfloat[H]{\includegraphics[width=\linewidth]{logo}}
\end{tabularx}

\caption{Many figures}\label{foo}
\end{figure}


% range of equations
http://tex.stackexchange.com/questions/18986/is-there-a-standard-way-to-refer-to-a-range-of-equations

% skip a section
\stepcounter{section}

Σχ.~\ref{<+Καλημέρα+>}<++>
$<++addhere++>\%$<++>


% put the figure in the exact place from top - useful for 2column docs
\begin{figure*}[t]
    \centering
    \includegraphics[trim=+3cm 0 0 -6cm,width=0.5\textwidth]{EMF_ex}
    \caption{Παραδείγματα χρήσης ΗΔΜ}
    \label{fig:lab2:EMF_ex}
\end{figure*}


% make a table fit a page - move it to the left
\begin{adjustwidth}{-2cm}{}
\end{adjustwidth}

% footnotes without markers
\newcommand\blfootnote[1]{%
  \begingroup
  \renewcommand\thefootnote{}\footnote{#1}%
  \addtocounter{footnote}{-1}%
  \endgroup
}

\usepackage{relsize} % for mathlarger

% remove any occurance of \textlatin command
:\\textlatin{\<\(\w\+\>\_s*\)*}
